% ======================================================================
% Chapter 21: Quantization of Spontaneously Broken Gauge Theories
% (book pages 731-773; PDF pages 754-803)
% ======================================================================
\chapter{自发破缺规范理论的量子化}

在第~20 章我们看到，当规范对称性被自发破缺时，规范玻色子获得质量。这一现象使我们得以构造弱相互作用的现实理论。然而，到目前为止，我们只是以一种简化的方式讨论了自发破缺规范理论。为了隔离物理自由度，我们使用了进入\emph{幺正规范}的手段。但是，在这个规范中微扰论的规则是什么，或者当我们计算费曼图时幺正规范约束如何得以保持，这些都完全不清楚。我们还看到，被吸收进有质量规范玻色子的戈德斯通玻色子在关于这些理论的形式论证中扮演着重要角色，因此我们希望以一种并不从一开始就消除这些粒子的规范来量子化这些理论。

在本章中，我们将通过用 Faddeev-Popov 方法对具有自发破缺规范对称性的理论做形式化的规范固定，来处理这些问题。我们将定义一族称为\emph{$R_\xi$ 规范}的规范，其中几乎全部包含原始自发对称性破缺的戈德斯通玻色子。这些粒子抵消形式体系中其他非物理粒子的效应，以保持理论的幺正性。这些抵消是我们第~16 章中看到的规范与鬼场自由度之间抵消的一个更精细的版本。然而，我们将在第~21.2 节看到，理论不会忘记它含有戈德斯通玻色子，并且在某些情形下，没有规范耦合的理论中戈德斯通玻色子的性质可以延续到含有有质量规范玻色子的理论中。

最后，在定义了微扰论并阐明了戈德斯通玻色子在自发破缺规范理论中的作用之后，我们将进行弱相互作用理论中一些感兴趣的显式圈计算。在这里我们将看到第~11 章思想的运用，即具有自发破缺对称性的理论可以用对称拉格朗日量的抵消项来重整化。在第~21.3 节中，我们将通过一些例子表明，这一结果对规范理论同样适用，并且它赋予弱相互作用规范理论以实质性的预言能力。

\section{$R_\xi$ 规范}

在我们对弱相互作用低能有效拉格朗日量的讨论中，我们在式~(20.89) 中提出了有质量规范玻色子传播子的如下表达式：
\begin{equation}
  \langle A^\mu(p)A^\nu(-p)\rangle = \frac{i}{p^2 - m^2}\ \Bigl(
    -g^{\mu\nu} + \frac{p^\mu p^\nu}{m^2} \Bigr) .
  \label{eq:21.1}
\end{equation}
这个表达式是一个自然的初步猜测，是对 Feynman-'t Hooft 规范的推广。然而，它在若干方面不能令人满意。

这些缺陷中最重要的一点涉及规范玻色子极化态的处理。传播子 \eqref{eq:21.1} 包含四个分量，分别对应横向、纵向和类时极化。我们在第~5 章和第~16 章看到，对于无质量规范玻色子，非物理的纵向和类时分量在计算中会抵消。然而，对于有质量规范玻色子，纵向极化态对应于一个真实的物理粒子；我们不希望它被抵消。表达式 \eqref{eq:21.1} 没有考虑这一变化。

\subsection{一个阿贝尔例子}

为了理解这一问题和具有自发破缺对称性的规范理论中出现的其他形式问题，我们需要仔细地重新做这些理论的 Faddeev-Popov 量子化。首先，我们将量子化式~\eqref{eq:20.1} 中引入的自发破缺阿贝尔规范理论：
\begin{equation}
  \mathcal{L} = -\frac{1}{4}(F_{\mu\nu})^2 + |D_\mu\phi|^2 - V(\phi) ,
  \label{eq:21.2}
\end{equation}
其中 $D_\mu = \partial_\mu + ieA_\mu$。这里 $\phi(x)$ 是复标量场。然而，把 $\phi$ 用其实分量写出来分析这个模型将是最方便的，
\begin{equation}
  \phi = \frac{1}{\sqrt{2}}(\phi_1 + i\phi_2) .
  \label{eq:21.3}
\end{equation}
于是无穷小局域对称变换为
\begin{equation}
  \delta\phi_1 = -\alpha(x)\phi_2,\qquad \delta\phi_2 = \alpha(x)\phi_1,\qquad
  \delta A_\mu = -\frac{1}{e}\partial_\mu\alpha .
  \label{eq:21.4}
\end{equation}
让我们假设 $V(\phi)$ 迫使标量场获得真空期望值：$\langle\phi\rangle = v$。然后我们应当通过一个平移来更换变量：
\begin{equation}
  \phi_1(x) = v + h(x);\qquad \phi_2 = \varphi .
  \label{eq:21.5}
\end{equation}
场 $\phi_2$ 即 $\varphi$ 是戈德斯通玻色子。拉格朗日量 \eqref{eq:21.2} 现在具有如下形式
\begin{equation}
  \mathcal{L} = -\frac{1}{4}(F_{\mu\nu})^2
    + \frac{1}{2}\bigl( \partial_\mu h - eA_\mu\varphi \bigr)^2
    + \frac{1}{2}\bigl( \partial_\mu\varphi + eA_\mu(v+h) \bigr)^2
    - V(\phi) .
  \label{eq:21.6}
\end{equation}
这个拉格朗日量仍然在一个精确的局域对称性下不变，
\begin{equation}
  \delta h = -\alpha\,\varphi,\qquad \delta\varphi = \alpha(v+h),\qquad
  \delta A_\mu = -\frac{1}{e}\partial_\mu\alpha .
  \label{eq:21.7}
\end{equation}
因此，为了定义对变量 $(\varphi, h, A^\mu)$ 的泛函积分，我们必须引入 Faddeev-Popov 规范固定。

从泛函积分出发
\begin{equation}
  Z = \int \mathcal{D}A\,\mathcal{D}h\,\mathcal{D}\varphi\ e^{i\int
    d^4x\ \mathcal{L}[A,h,\varphi]} ,
  \label{eq:21.8}
\end{equation}
我们可以像第~9.4 节那样引入一个规范固定约束。按照从式~(9.50) 到式~\eqref{eq:9.54} 的步骤，我们求得
\begin{equation}
  Z = C\cdot \int \mathcal{D}A\,\mathcal{D}h\,\mathcal{D}\varphi\ e^{i\int
    d^4x\ \mathcal{L}[A,h,\varphi]}\ \delta\bigl( G(A,h,\varphi) \bigr)\
    \det\Bigl( \frac{\delta G}{\delta\alpha} \Bigr) ,
  \label{eq:21.9}
\end{equation}
其中 $C$ 是一个正比于规范群体积的常量，$G(A,h,\varphi)$ 是规范固定条件。或者，我们可以像式~\eqref{eq:9.56} 的推导那样，把规范固定约束引入为 $\delta(G(x) - \omega(x))$ 并用高斯权重对 $\omega(x)$ 积分。这给出
\begin{equation}
  Z = C'\cdot \int \mathcal{D}A\,\mathcal{D}h\,\mathcal{D}\varphi\ \exp\Bigl[
    i\int d^4x\ \Bigl( \mathcal{L}[A,h,\varphi] - \frac{1}{2\xi}(G)^2
      \Bigr) \Bigr]\ \det\Bigl( \frac{\delta G}{\delta\alpha} \Bigr) .
  \label{eq:21.10}
\end{equation}
规范固定函数 $G$ 是任意的，但我们可以通过恰当地选择它来简化我们的形式体系。

规范固定函数一个特别方便的选择是
\begin{equation}
  G = \partial^\mu A_\mu - \xi e v\,\varphi .
  \label{eq:21.11}
\end{equation}
当我们构成 $G^2$ 时，$A^\mu$ 的二次项将为规范场作用量提供我们在式~\eqref{eq:9.58} 和 \eqref{eq:16.29} 的推导中看到的同样依赖于规范的附加项。此外，$A^\mu$ 与 $\varphi$ 之间的交叉项被设计成抵消来自 \eqref{eq:21.6} 第三项的 $\partial^\mu A_\mu\,\varphi$ 形式的二次项。采用这一选择，规范固定的拉格朗日量 $(\mathcal{L} - \frac{1}{2\xi}G^2)$ 的二次项为
\begin{equation}
  \begin{aligned}
    \mathcal{L}_{\mathrm{quad}} &= -\frac{1}{4}(F_{\mu\nu})^2
      - \frac{1}{2\xi}(\partial^\mu A_\mu)^2
      + \frac{1}{2}m_A^2 A_\mu A^\mu + \frac{1}{2}m_A^2\,\varphi^2
      + \frac{1}{2}m_h^2 h^2 .
  \end{aligned}
  \label{eq:21.12}
\end{equation}
$h$ 场的质量项来自 $V(\phi)$ 的展开，如同 \eqref{eq:20.6}。规范场的质量项来自希格斯机制，即来自 \eqref{eq:21.6} 的第三项。注意，这个形式体系还为戈德斯通玻色子 $\varphi$ 产生了一个质量：
\begin{equation}
  m_\varphi^2 = \xi(e v)^2 = \xi m_A^2 .
  \label{eq:21.13}
\end{equation}
这个质量依赖规范的事实是戈德斯通玻色子是一个虚构场的信号，它在物理过程中不会被产生。

为了完成 Faddeev-Popov 量子化步骤，我们必须推导鬼场的拉格朗日量。这个拉格朗日量依赖于 $G$ 的规范变化，它可以通过把 \eqref{eq:21.7} 代入 \eqref{eq:21.11} 来计算。我们求得
\begin{equation}
  \delta G = \frac{1}{e}\bigl( -\partial^2 - \xi e^2v^2 \bigr)\alpha
  = \frac{1}{e}\bigl( -\partial^2 - \xi m_A^2 \bigr)\alpha .
  \label{eq:21.14}
\end{equation}
这个算符的行列式可以通过引入一组具有如下拉格朗日量的 Faddeev-Popov 鬼场来计入，
\begin{equation}
  \mathcal{L}_{\mathrm{ghost}} = \bar{c}\bigl( -\partial^2 - \xi m_A^2 \bigr)c ,
  \label{eq:21.15}
\end{equation}
其中 $m_A = ev$，如同式~\eqref{eq:21.13}。由于这是阿贝尔规范理论，鬼场不直接耦合到规范场。然而，它的确耦合到物理的希格斯场，因此不能像在 QED 中那样被完全忽略。

从 $A^\mu$、$h$、$\varphi$ 和鬼场的拉格朗日量中的二次项，我们可以容易地求得这些场的传播子。全部四个传播子都显示在图~\ref{fig:21.1} 中。唯一复杂的情形是规范场。\eqref{eq:21.12} 中涉及 $A^\mu$ 的项包含一个算符，其 Fourier 变换为
\begin{equation}
  g^{\mu\nu}k^2 - \Bigl( 1 - \frac{1}{\xi} \Bigr)k^\mu k^\nu
    - m_A^2 g^{\mu\nu} .
  \label{eq:21.16}
\end{equation}
这个矩阵的逆给出 $A^\mu$ 场传播子：
\begin{equation}
  \frac{-i}{k^2 - m_A^2}\ \Bigl( g^{\mu\nu} - (1-\xi)\frac{k^\mu k^\nu}
    {k^2 - \xi m_A^2} \Bigr) .
  \label{eq:21.17}
\end{equation}
注意，$A$ 场的横向分量和希格斯场的分量 $h$ 获得了我们在第~20.1 节中求得的质量 $m_A$ 和 $m_h$。$A$ 的非物理分量、戈德斯通玻色子和鬼场都获得相同的依赖规范的质量 $\sqrt{\xi}m_A$。

\begin{figure}[t]
  \centering
  \includegraphics[width=0.55\textwidth]{fig21.1.png}
  \caption{自发破缺对称性的阿贝尔模型中规范场、希格斯场和鬼场的传播子。}
  \label{fig:21.1}
\end{figure}

\subsection{微扰论中的 $\xi$ 依赖性}

由于参数 $\xi$ 只是在规范固定中被引入，我们预期它会从规范不变算符期望值和 $S$ 矩阵元的所有计算中消去。利用规范固定拉格朗日量的 BRST 对称性，这种抵消可以在微扰论的所有阶被证明。%
\footnote{参见，例如，Taylor (1976)。} 然而，在这里我们将只是在一个简单例子中说明 $\xi$ 的抵消。

\begin{figure}[t]
  \centering
  \includegraphics[width=0.55\textwidth]{fig21.2.png}
  \caption{具有自发对称性破缺的阿贝尔模型中对费米子-费米子散射有领头阶贡献的图。}
  \label{fig:21.2}
\end{figure}

考虑通过一个手征相互作用把费米子耦合到自发破缺规范理论：
\begin{equation}
  \mathcal{L}_f = \bar{\psi}_L(i\slashed{D})\psi_L
    + \bar{\psi}_R(i\slashed{\partial})\psi_R
    - \lambda\bigl( \bar{\psi}_R\phi\,\psi_L + \bar{\psi}_L\phi^{*}\psi_R
      \bigr) ,
  \label{eq:21.18}
\end{equation}
其中 $D_\mu = \partial_\mu + ieA_\mu$，同前。这是费米子与弱相互作用规范理论耦合的一个简化了的阿贝尔版本。费米子 $\psi$ 从自发对称性破缺中获得质量
\begin{equation}
  m_f = \frac{\lambda v}{\sqrt{2}}
  \label{eq:21.19}
\end{equation}
（这个理论具有一个轴矢量反常，它会使圈计算不自洽，但我们只会在树图级别分析它。）

在这个理论中，对费米子-费米子散射有领头阶贡献的图是图~\ref{fig:21.2} 所示的那些图。注意，必须包括非物理粒子 $\varphi$ 交换的贡献，因为这个粒子出现在费曼规则中。鬼场直到一圈级别才出现在这个过程中。由于物理的希格斯粒子 $h$ 的传播子与 $\xi$ 无关，$\xi$ 依赖性的抵消必须发生在 $A^\mu$ 的横向与纵向分量和戈德斯通玻色子 $\varphi$ 之间。

显式地进行这种抵消，人们发现纵向规范玻色子传播子的贡献与戈德斯通玻色子交换相结合，给出规范不变的结果。这验证了 $R_\xi$ 规范给出与 $\xi$ 无关的结果，正如规范不变性所要求的那样。

\subsection{Feynman-'t Hooft 规范与幺正规范}

在 $R_\xi$ 规范中，选择 $\xi = 1$ 特别方便。在这个被称为\emph{Feynman-'t Hooft 规范}的规范中，规范玻色子和戈德斯通玻色子的传播子具有简单形式
\begin{equation}
  \frac{-ig^{\mu\nu}}{k^2 - m_A^2},\qquad
  \frac{i}{k^2 - m_A^2} .
\end{equation}
这个规范仍被称为 Feynman-'t Hooft 规范，是进行一般高阶计算最方便的规范。

对 $\xi$ 的任何有限值，规范玻色子和戈德斯通玻色子的传播子都按 $1/k^2$ 衰减，因而服从第~10.1 节的一般幂次计数分析。由此推出，在任一这些规范中，微扰论都将是可重整化的，其意义是发散被有限个抵消项消除。此外，第~11.6 节的分析告诉我们，所需的唯一抵消项是在理论的原始整体对称性下对称的那些项。然而，我们应当对我们的重整化步骤要求进一步的条件：我们应当坚持抵消项保持局域规范不变性，特别是保持 $S$ 矩阵元和规范不变算符的矩阵元与 $\xi$ 无关这一性质。这一结果由 't Hooft 和 Veltman 以及 Lee 和 Zinn-Justin 在微扰论的所有阶证明。%
\footnote{G.~'t~Hooft 和 M.~J.~G.~Veltman，\emph{Nucl.\ Phys.\ B} \textbf{50}，318 (1972)；B.~W.~Lee 和 J.~Zinn-Justin，\emph{Phys.\ Rev.\ D} \textbf{5}，3121、3137、3155 (1972)，\textbf{D7}，1049 (1973)。} 因此，在由 $\xi$ 的任何有限值定义的规范中，原则上我们可以直截了当地把物理量计算到任意阶。由 $\xi$ 的可能值定义的这些规范被称为\emph{可重整化}规范，即 $R_\xi$ 规范。

通过取 $R_\xi$ 规范的极限 $\xi\to\infty$，我们找到一个具有非常不同的简化特征的规范。在这个极限下，质量正比于 $\sqrt{\xi}$ 的非物理自由度从理论中消失。规范玻色子和戈德斯通玻色子的传播子变为：
\begin{equation}
  \frac{-i}{k^2 - m_A^2}\ \Bigl( g^{\mu\nu} - \frac{k^\mu k^\nu}{m_A^2}
    \Bigr),\qquad \frac{i}{k^2 - \xi m_A^2}\ \to 0 .
\end{equation}
规范玻色子传播子恰好包含三个类空极化态。在这个规范中，费曼图的唯一奇点对应于物理中间态的传播。因此，$S$ 矩阵的幺正性由 Cutkosky 规则推出，如同第~7.3 节中考虑的整体对称理论那样，无需担心非物理态的抵消。%
\footnote{在第 16.4 节更精细的语言中，$S$ 矩阵的幺正性所必需的关键恒等式 \eqref{eq:16.54} 是显然成立的。} 因此，$R_\xi$ 规范的 $\xi\to\infty$ 极限给出了式~\eqref{eq:20.12} 中引入的\emph{幺正}（或 $U$）规范的量子力学实现。

直接证明 $U$ 规范中的可重整化性并不直截了当。在这个规范中，规范玻色子传播子在大的 $k$ 处衰减得比 $1/k^2$ 更慢。这预示了圈图计算中的麻烦。事实上，典型地，个别的圈图会随 $\xi\to\infty$ 发散为 $\log\xi$ 或更差。尽管如此，$S$ 矩阵的规范不变性意味着这些发散必须在贡献给给定过程的所有图之和中抵消，使得这个和在 $\xi\to\infty$ 时有光滑极限。我们用一个规范证明自发破缺规范理论的可重整化性、用另一个规范证明其幺正性，这一事实在原则上并无困难。事实上，这种论证方法自然地利用了理论的基本对称性。

\subsection{非阿贝尔分析}

既然我们已经透彻地考察了阿贝尔规范理论的 $R_\xi$ 规范，我们准备推广到非阿贝尔情形。做到完全一般并无困难，因此让我们考虑一个规范群为 $G$、由标量场的真空期望值自发破缺的 Yang-Mills 规范理论。

我们将以式~\eqref{eq:20.13} 之后对这个系统的经典分析为基础。如同那个分析，把标量场写成实值场的多重态 $\phi_i$ 将是最方便的。于是 $\phi_i$ 的规范变换具有如下形式
\begin{equation}
  \delta\phi_i = -\alpha^a(x)\,T^a_{ij}\phi_j
  \label{eq:21.30}
\end{equation}
其中 $T^a$ 是 $G$ 的实的、反对称的表示矩阵。类似地，规范场的变换为
\begin{equation}
  \delta A^a_\mu = \partial_\mu\alpha^a - f^{abc}\alpha^b A^c_\mu
  = (D_\mu\alpha)^a .
  \label{eq:21.31}
\end{equation}
（如果规范群不是单群，耦合常数 $g$ 不必对每个 $a$ 都相同。）在这些规范变换下不变的拉格朗日量是
\begin{equation}
  \mathcal{L} = -\frac{1}{4}(F^a_{\mu\nu})^2 + \frac{1}{2}|D_\mu\phi|^2
    - V(\phi) ,
  \label{eq:21.32}
\end{equation}
其中
\begin{equation}
  D_\mu\phi_i = \partial_\mu\phi_i + gA^{a\mu}\,T^a_{ij}\phi_j .
  \label{eq:21.33}
\end{equation}
假设势 $V(\phi_i)$ 在 $\phi$ 的某些分量获得真空期望值的点处取极小。如同 \eqref{eq:20.16}，定义
\begin{equation}
  \langle\phi_i\rangle = \phi_{0i} .
  \label{eq:21.34}
\end{equation}
我们将围绕这个值展开 $\phi_i$：
\begin{equation}
  \phi_i(x) = \phi_{0i} + \chi_i(x) .
  \label{eq:21.35}
\end{equation}
把 $\chi_i$ 的值空间分成两个子空间将是很方便的。矢量 $T^a\phi_0$ 对应于 $\phi$ 真空期望值的对称变换。沿这些方向的场涨落就是戈德斯通玻色子。设 $\{\eta_i^k\}$ 是这个子空间的正交归一基；于是单位矢量 $\eta_i^k$ 与戈德斯通玻色子一一对应。与所有矢量 $T^a\phi_0$ 正交的场涨落对应于自发破缺规范理论的（有质量的）物理标量场。

在随后的讨论中，矢量 $T^a\phi_0$ 将扮演重要角色。因此我们应当回忆在式~\eqref{eq:20.51} 中引入的这些矢量的记号：
\begin{equation}
  F^{ai} = T^a_{ij}\phi_{0j} .
  \label{eq:21.36}
\end{equation}
矩阵 $F^{ai}$ 一般不是方阵；它对每个规范生成元有一行，对 $\phi$ 的每个分量有一列。然而，它的许多元素为零。它的非零元素连接自发破缺的规范生成元和戈德斯通玻色子。在式~\eqref{eq:20.56} 中，我们证明了通过希格斯机制产生的规范玻色子质量可以写成
\begin{equation}
  m^2_{ab} = g^2\,F^{aj}F^{bj} .
  \label{eq:21.37}
\end{equation}

为了给出矩阵 $F^{aj}$ 的一个具体例子，让我们在 GWS 电弱理论中计算它。遵循式~\eqref{eq:20.14} 中引入的约定，我们应当把 GWS 模型的希格斯场用四个实标量场重写。一个方便的参数化是
\begin{equation}
  \phi = \frac{1}{\sqrt{2}}\begin{pmatrix}
    \phi_1 + i\phi_2\\
    v + h + i\phi_3
  \end{pmatrix} .
  \label{eq:21.38}
\end{equation}
场 $\phi^1, \phi^2, \phi^3$ 是戈德斯通玻色子，而 $h$ 是有质量的希格斯玻色子。真空态就是
\begin{equation}
  \phi_0 = \frac{1}{\sqrt{2}}\begin{pmatrix}0\\ v\end{pmatrix} .
\end{equation}
实的表示矩阵是
\begin{equation}
  T^a = -i\tau^a = -i\frac{\sigma^a}{2},\qquad
  T^Y = -iY = -i\frac{1}{2} .
\end{equation}
于是简单计算表明，例如，$T^1\phi_0$ 等于 $v/2$ 乘以 1 方向上的单位矢量。填入 $F^{ai}$ 中 $a = 1, 2, 3, Y$ 和 $i = 1, 2, 3$ 的其余分量，我们求得
\begin{equation}
  F^{ai} = \frac{v}{2}\ \begin{pmatrix}
    1 & 0 & 0\\
    0 & 1 & 0\\
    0 & 0 & 1\\
    0 & 0 & 0
  \end{pmatrix},
  \label{eq:21.39}
\end{equation}
我们不需要包括 $F^{ai}$ 沿物理希格斯场 $h$ 方向的分量；矢量 $T^a\phi_0$ 都与此方向正交。

如果我们把 \eqref{eq:21.35} 作为变量更换代入 \eqref{eq:21.32}，我们求得，对拉格朗日量中的二次项，
\begin{equation}
  \mathcal{L}_2 = -\frac{1}{4}A^{a\mu}\bigl( -g^{\mu\nu}\partial^2
    + \partial^\mu\partial^\nu \bigr)A^a_\nu + \frac{1}{2}(\partial_\mu\chi)^2
    + g\,\partial_\mu\chi_i\,A^{a\mu}F^{ai}
    + \frac{1}{2}(m^2_A)_{ab}A^a_\mu A^{b\mu} ,
  \label{eq:21.40}
\end{equation}
其中 $(m^2_A)_{ab}$ 是规范玻色子质量矩阵 \eqref{eq:21.37}，且
\begin{equation}
  M^2_{ij} = \frac{\partial^2V}{\partial\phi_i\partial\phi_j}
    \Big|_{\phi_0} .
  \label{eq:21.41}
\end{equation}
我们在式~\eqref{eq:11.13} 中证明了
\begin{equation}
  \eta_i^k M^2_{ij} = 0
  \label{eq:21.42}
\end{equation}
对由 $T^a\phi_0$ 张成的子空间中的一切可能方向 $\eta_i^k$ 成立，因此戈德斯通玻色子是无质量的。

为了研究这个系统的量子理论，我们从泛函积分出发
\begin{equation}
  Z = \int \mathcal{D}A\,\mathcal{D}\chi\ e^{i\int d^4x\ \mathcal{L}[A,\chi]} .
  \label{eq:21.43}
\end{equation}
利用 Faddeev-Popov 规范固定步骤，我们把这个积分类比于 \eqref{eq:21.10} 定义为
\begin{equation}
  Z = C'''\cdot \int \mathcal{D}A\,\mathcal{D}\chi\ \exp\Bigl[ i\int d^4x\
    \Bigl( \mathcal{L}[A,\chi] - \frac{1}{2\xi}(G)^2 \Bigr) \Bigr]\
    \det\Bigl( \frac{\delta G}{\delta\alpha} \Bigr) ,
  \label{eq:21.44}
\end{equation}
对任意的规范固定函数 $G(A,\chi)$。$R_\xi$ 规范由如下选择定义
\begin{equation}
  G^a = \partial^\mu A^a_\mu - \xi g F^{ai}\chi_i .
  \label{eq:21.45}
\end{equation}
注意，$G$ 只涉及位于戈德斯通玻色子子空间中的 $\chi$ 分量。

规范固定项把如下一组二次项加到拉格朗日量中：
\begin{equation}
  \Bigl( -\frac{1}{2\xi}G^2 \Bigr) = \frac{1}{2\xi}A^{a\mu}\bigl(
    \partial^\mu\partial^\nu \bigr)A^a_\nu + g\,\partial^\mu A^{a\mu}F^{ai}\chi_i
    - \frac{\xi g^2}{2}\bigl[ F^{ai}\chi_i \bigr]^2 .
  \label{eq:21.46}
\end{equation}
混合 $A^a_\mu$ 和 $\chi_i$ 的项被安排在 \eqref{eq:21.40} 与 \eqref{eq:21.46} 之间抵消。规范场和戈德斯通玻色子场的最终二次拉格朗日量是
\begin{equation}
  \mathcal{L}_2 = -\frac{1}{2}A^{a\mu}\Bigl[ \Bigl( -g_{\mu\nu}\partial^2
    + \Bigl( 1 - \frac{1}{\xi} \Bigr)\partial_\mu\partial_\nu
    - g^2F^{ai}F^{bi}g_{\mu\nu} \Bigr)A^{b\nu} \Bigr]
    - \frac{1}{2}\xi g^2(F^{ai}\chi_i)^2 .
  \label{eq:21.47}
\end{equation}
这个拉格朗日量中规范玻色子和戈德斯通玻色子的质量矩阵彼此密切相关。规范玻色子质量矩阵是
\begin{equation}
  (m^2_A)_{ab} = g^2F^{ai}F^{bi} = g^2(FF^T)_{ab} .
  \label{eq:21.48}
\end{equation}
在 $R_\xi$ 规范中，规范玻色子的类时分量获得质量
\begin{equation}
  \xi m^2_A = \xi g^2(FF^T)_{ab} .
  \label{eq:21.49}
\end{equation}
同时，戈德斯通玻色子获得质量矩阵
\begin{equation}
  m^2_\varphi = \xi g^2(F^TF)_{ij} .
  \label{eq:21.50}
\end{equation}
两个矩阵 \eqref{eq:21.49} 和 \eqref{eq:21.50} 具有不同数目的零本征值，但它们的非零本征值一一对应。这正是希格斯机制在有质量规范玻色子与它们为获得质量而吸收的戈德斯通玻色子之间所诱导的对应。

最后，我们必须构造鬼场拉格朗日量。这从规范固定项 $G^a$ 的规范变化求得。把 \eqref{eq:21.30} 和 \eqref{eq:21.31} 代入 \eqref{eq:21.45}，我们求得
\begin{equation}
  \delta G^a = \Bigl( -\partial^2\delta^{ab} + \xi g^2\,(T^a\phi_0)\cdot
    T^b(\phi_0 + \chi) \Bigr)\alpha^b .
  \label{eq:21.51}
\end{equation}
因此，鬼场拉格朗日量是
\begin{equation}
  \mathcal{L}_{\mathrm{ghost}} = \bar{c}^a\Bigl[ -\partial^2\delta^{ab}
    - \xi g^2\,(T^a\phi_0)\cdot T^b(\phi_0 + \chi) \Bigr]c^b .
  \label{eq:21.52}
\end{equation}
注意，鬼场具有与规范玻色子的非物理分量完全相同的质量矩阵 \eqref{eq:21.49}。这个拉格朗日量还同时包含鬼场到规范场的熟悉耦合，以及我们在阿贝尔情形 \eqref{eq:21.15} 中求得的对物理希格斯场的耦合。

我们现在已经计算了 $R_\xi$ 规范中规范场、标量场和鬼场的动能项。把这些结果转化为这些场的传播子计算是直截了当的；这些计算与阿贝尔情形完全相同。我们求得三个传播子的标准形式，其规范玻色子和戈德斯通玻色子质量由矩阵 \eqref{eq:21.48}--\eqref{eq:21.50} 给出。这就完成了 $R_\xi$ 规范中自发破缺规范理论的量子化。

\section{戈德斯通玻色子等价定理}

$R_\xi$ 规范中自发破缺规范理论的量子化导致一个卓越而重要的结果，即\emph{戈德斯通玻色子等价定理}：在高能下，发射或吸收纵向极化规范玻色子的振幅变得等于相应戈德斯通玻色子的振幅。这个定理对弱相互作用的物理有重要应用，因为 $W$ 和 $Z$ 玻色子的纵向分量在高能下主导它们的产生和衰变。

为准备对纵向矢量玻色子的讨论，我们需要一些简单的运动学。一个静止的矢量玻色子具有动量 $k^\mu = (m, 0, 0, 0)$ 和一个极化矢量，后者是三个正交单位矢量的线性组合
\begin{equation}
  (0, 1, 0, 0),\qquad (0, 0, 1, 0),\qquad (0, 0, 0, 1) .
  \label{eq:21.57}
\end{equation}
如果我们沿 3 轴对这个粒子作 boost，其动量变为
\begin{equation}
  k^\mu = (E_k, 0, 0, k) .
\end{equation}
三个可能的极化矢量现在是满足如下条件的三个单位矢量
\begin{equation}
  \epsilon^2 = 0,\qquad \epsilon^{*}\cdot\epsilon = -1 .
  \label{eq:21.58}
\end{equation}
其中两个是 \eqref{eq:21.57} 中的前两个矢量；它们给出横向极化。满足 \eqref{eq:21.58} 的第三个矢量是纵向极化矢量
\begin{equation}
  \epsilon^\mu_L(k) = \Bigl( \frac{k}{m}, 0, 0, \frac{E_k}{m} \Bigr) ,
  \label{eq:21.59}
\end{equation}
它是 \eqref{eq:21.57} 中第三个矢量的 boost。\eqref{eq:21.59} 的一个重要且多少有些反直觉的特征是，当 $k$ 变大时它变得越来越平行于 $k^\mu$。事实上，逐分量地，
\begin{equation}
  \epsilon^\mu_L(k) = \frac{k^\mu}{m} + \mathcal{O}(m/E_k)
  \label{eq:21.60}
\end{equation}
当 $k\to\infty$ 时。由于 $k^\mu$ 的分量随 $k$ 增长，这一陈述与要求 $\epsilon^{*}\cdot k = 0$ 而 $k\cdot k = m^2$ 是一致的。

记住这个运动学情形，让我们分析壳上态之间规范流矩阵元满足的 Ward 恒等式。最简便的是在 Lorentz 规范（$\xi = 0$）中处理，其中规范固定项 \eqref{eq:21.45} 不涉及戈德斯通玻色子场。于是 Ward 恒等式可以写成如下形式：
\begin{equation}
  k_\mu\langle \beta|\ J^{\mu a}(k)\ |\alpha\rangle = 0 ,
  \label{eq:21.64}
\end{equation}
在壳上态之间成立。把矩阵元写成两块之和---第一，流可以直接耦合进一个单粒子不可约顶角函数 $\Gamma^\mu(k)$；第二，对自发破缺规范理论，还有一个额外的非单粒子不可约项，其中流产生一个戈德斯通玻色子，正是这个粒子通过一个 1PI 顶角 $\Gamma(k)$ 耦合到外态---我们求得
\begin{equation}
  k_\mu\Gamma^\mu(k) + k_\mu\,(igF)\ \frac{1}{k^2}\ \Gamma(k) = 0 .
  \label{eq:21.65}
\end{equation}
让我们把联系规范流与戈德斯通玻色子态的关系写成
\begin{equation}
  \langle 0|J^\mu|\pi(k)\rangle = -iF k^\mu ,
  \label{eq:21.62}
\end{equation}
如同式~\eqref{eq:20.46}。于是导致式~\eqref{eq:20.56} 的论证告诉我们，规范玻色子质量由下式给出
\begin{equation}
  m = gF ,
  \label{eq:21.63}
\end{equation}
其中 $g$ 是规范玻色子耦合常数。借助这些等同，我们可以把 Ward 恒等式 \eqref{eq:21.64} 写成
\begin{equation}
  k_\mu\Gamma^\mu(k) = m\,\Gamma(k) .
  \label{eq:21.66}
\end{equation}
现在在规范玻色子动量大时使用这个方程。由于规范玻色子顶角是单粒子不可约的，顶角内部传播子的动量一般与 $k^\mu$ 不共线。于是，根据 \eqref{eq:21.60}，我们可以用纵向极化矢量替换 $k^\mu/m$。注意，在 \eqref{eq:21.65} 的第二项中这既不允许（而且也不需要）。我们的最终结果是
\begin{equation}
  \epsilon^\mu_L(k)\,\Gamma_\mu(k) = \Gamma(k) ,
  \label{eq:21.67}
\end{equation}
当 $k\to\infty$ 时成立，误差为 $m^2/k^2$ 量级。也就是说，在高能极限下，纵向规范玻色子的耦合变得精确地等于其关联戈德斯通玻色子的耦合。

等价定理可以用另一种方式导出，利用我们在式~(21.26) 之后讨论的自发破缺规范理论中物理态的计数。在上一节我们看到，至少在树图级别，自发破缺规范理论中幺正性通过产生类时极化规范玻色子的图与产生戈德斯通玻色子的图之间的抵消来维持。

在 Feynman-'t Hooft 规范中情形最为清晰。在那里，规范玻色子传播子的分子是 $-g^{\mu\nu}$。我们可以用极化矢量把它写成
\begin{equation}
  -g^{\mu\nu} = \sum_{\epsilon} \epsilon^{\mu*}(k)\epsilon^\nu(k)
    - \frac{k^\mu k^\nu}{m^2} .
  \label{eq:21.68}
\end{equation}
最后一项是来自非物理类时极化态的贡献。$S$ 矩阵的幺正性要求，当穿过一个图的 Cutkosky 割线把规范玻色子传播子放到壳上时，这一部分的贡献必须被穿过戈德斯通玻色子线的 Cutkosky 割线抵消。所需的抵消是
\begin{equation}
  \Bigl| \frac{k^\mu}{m}\Gamma_\mu(k) \Bigr|^2 + |\Gamma(k)|^2 = 0 ,
  \label{eq:21.69}
\end{equation}
或者，用图来表示，
\begin{equation}
  \text{(timelike gauge boson cut)} + \text{(Goldstone boson cut)} = 0 .
\end{equation}
再次地，由于 $\Gamma^\mu(k)$ 是单粒子不可约顶角，对高能规范玻色子我们可以用 \eqref{eq:21.60} 把 $(k^\mu/m)$ 替换为纵向极化矢量 $\epsilon^\mu_L(k)$。于是 \eqref{eq:21.69} 变成正是 \eqref{eq:21.67} 的平方。

通过这些形式论证，我们至少可以在单规范玻色子发射过程的树图级别看到，等价定理必然成立。然而，在有趣物理过程的显式计算中看到等价定理起作用要更有启发性。我们现在将用两个例子说明它的影响。%
\footnote{关于等价定理的一般证明，包括有质量矢量玻色子的多重吸收和发射情形，参见 M.~S.~Chanowitz 和 M.~K.~Gaillard，\emph{Nucl.\ Phys.\ B} \textbf{261}，379 (1985)。}

\subsection{顶夸克衰变}

第一个例子是顶夸克的弱衰变。这个电荷为 $+2/3$ 的夸克足够重，可以通过 $t\to W^+ + b$ 衰变成一个实的 $W^+$。这个衰变的图由图~\ref{fig:21.4} 所示的简单规范顶角给出。

\begin{figure}[t]
  \centering
  \includegraphics[width=0.45\textwidth]{fig21.4.png}
  \caption{$t$ 夸克衰变为 $W^+ + b$。}
  \label{fig:21.4}
\end{figure}

让我们首先尝试猜测顶夸克衰变宽度的大小。矩阵元的平方将包含一个 $g^2$ 因子，乘以某个具有质量量纲的表达式。由于如果顶夸克质量重，宽度应当大，一个初步猜测可能是
\begin{equation}
  \Gamma \sim \frac{g^2}{4\pi}\,m_t .
  \label{eq:21.70}
\end{equation}
然而，正确的表达式结果证明被增强了 $(m_t/m_W)^2$ 倍。

这个衰变的振幅可以从式~\eqref{eq:20.80} 读出：
\begin{equation}
  i\mathcal{M} = \frac{ig}{\sqrt{2}}\ \bar{u}(q)\gamma^\mu\Bigl(
    \frac{1-\gamma^5}{2} \Bigr)u(p)\,\epsilon^{*}_\mu(k) .
  \label{eq:21.71}
\end{equation}
（我们把相关的 CKM 因子设为 1。）我们现在把这个振幅转化为顶夸克衰变率的表达式。为简单起见，在这个计算中我们将忽略 $b$ 夸克的质量。

按照我们的标准方法把 \eqref{eq:21.71} 中的振幅平方，然后对初态自旋求平均并对末态自旋求和，我们求得
\begin{equation}
  \sum_{\text{spins}}|\mathcal{M}|^2 = \frac{g^2}{2}\ \mathrm{tr}\Bigl[
    \slashed{q}\gamma^\mu\slashed{p}\gamma^\nu\ \frac{1-\gamma^5}{2}
    \Bigr]\ \sum_{\text{polarizations}}\ \epsilon^{*}_\mu(k)\epsilon_\nu(k)
  .
  \label{eq:21.72}
\end{equation}
通过插入极化求和的表达式，我们可以显式地对物理规范玻色子极化求和。这给出
\begin{equation}
  \sum_{\text{spins}}|\mathcal{M}|^2 = \frac{g^2}{2}\ \Bigl[
    4\Bigl( -g^{\mu\nu} + \frac{k^\mu k^\nu}{m_W^2} \Bigr)\ \cdots
    \Bigr] ,
  \label{eq:21.73}
\end{equation}
其中正比于 $k^\mu k^\nu/m_W^2$ 的纵向极化贡献给出主导项。计算迹，人们发现纵向贡献被增强了 $(m_t/m_W)^2$ 倍，正如等价定理所要求。顶夸克的衰变宽度是
\begin{equation}
  \Gamma(t\to W^+b) = \frac{g^2}{64\pi}\ \frac{m_t^3}{m_W^2}\ \cdots ,
  \label{eq:21.74}
\end{equation}
它比朴素的估计 \eqref{eq:21.70} 大得多。这是戈德斯通玻色子等价定理的直接例证：末态中的纵向 $W$ 玻色子的行为像相应的戈德斯通玻色子，而振幅被大的 $t$ 夸克质量增强。

\subsection{$e^{+}e^{-}\to W^{+}W^{-}$}

第二个例子是弱相互作用规范理论中的过程 $e^{+}e^{-}\to W^{+}W^{-}$。相关图显示在图~\ref{fig:21.7} 中。这个过程很有意思，因为根据等价定理，两个纵向 $W$ 玻色子产生的截面在高能下必然等于两个戈德斯通玻色子产生的截面。

\begin{figure}[t]
  \centering
  \includegraphics[width=0.6\textwidth]{fig21.7.png}
  \caption{弱相互作用规范理论中对 $e^{+}e^{-}\to W^{+}W^{-}$ 有贡献的图。}
  \label{fig:21.7}
\end{figure}

让我们首先计算一对带电标量（戈德斯通玻色子）产生的振幅。从图~\ref{fig:21.8} 组装顶角，

\begin{figure}[t]
  \centering
  \includegraphics[width=0.6\textwidth]{fig21.8.png}
  \caption{弱相互作用规范理论中电子和标量耦合到光子和 $Z$ 玻色子的费曼规则。}
  \label{fig:21.8}
\end{figure}

我们求得，对任一螺旋度的电子，通过虚光子湮灭为标量的振幅是
\begin{equation}
  i\mathcal{M}(e\bar{e}\to\gamma^{*}\to\phi^+\phi^-)
  = ie^2\ \bar{v}(p')\gamma^\mu u(p)\ \frac{1}{q^2}\ (k_+ - k_-)_\mu ,
  \label{eq:21.92}
\end{equation}
其中 $k_+$、$k_-$ 是标量的动量，$q = k_+ + k_-$。通过虚 $Z^0$ 湮灭的相应振幅依赖于 $e^{+}e^{-}$ 的螺旋度。把这些贡献加到前面的表达式上，我们求得
\begin{equation}
  \begin{aligned}
    i\mathcal{M}(e^-_Le^+_R \to \phi^+\phi^-) &= ie^2\ \bar{v}_R\gamma^\mu
      u_L\ \Bigl[ \frac{1}{q^2} + \frac{\bigl( \tfrac{1}{2} - \sin^2\theta_W
        \bigr)^2}{\sin^2\theta_W\cos^2\theta_W}\ \frac{1}{q^2 - m_Z^2}
      \Bigr](k_+ - k_-)_\mu ,\\
    i\mathcal{M}(e^-_Re^+_L \to \phi^+\phi^-) &= ie^2\ \bar{v}_L\gamma^\mu
      u_R\ \Bigl[ \frac{1}{q^2} + \frac{(\sin^2\theta_W)^2}{\sin^2\theta_W
        \cos^2\theta_W}\ \frac{1}{q^2 - m_Z^2} \Bigr](k_+ - k_-)_\mu .
  \end{aligned}
  \label{eq:21.93}
\end{equation}
注意，在高能极限下，右旋电子湮灭的振幅约化为
\begin{equation}
  i\mathcal{M}(e^-_Re^+_L \to \phi^+\phi^-)
  = ie^2\ \bar{v}_L\gamma^\mu u_R\ \frac{1}{q^2}\,(k_+ - k_-)_\mu ,
  \label{eq:21.94}
\end{equation}
它正是 $Y = -1$ 的 $e_R$ 通过耦合常数为 $g' = e/\cos\theta_W$ 的 $U(1)$ 规范玻色子 $B_\mu$ 耦合到 $Y = 1/2$ 的 $\phi^+$ 的振幅。这个表达式反映了 $\phi^+$ 与 $SU(2)$ 规范玻色子没有直接耦合这一事实。类似地，左旋电子的振幅趋于
\begin{equation}
  i\mathcal{M}(e^-_Le^+_R \to \phi^+\phi^-)
  = ie^2\ \bar{v}_R\gamma^\mu u_L\ \frac{1}{4\cos^2\theta_W\sin^2\theta_W}\
    \frac{1}{q^2}\,(k_+ - k_-)_\mu ,
  \label{eq:21.95}
\end{equation}
在高能极限下成立。它具有 $B^\mu$ 和 $A^\mu$ 交换振幅相干叠加的结构。正像我们在第~11 章看到的那样，具有自发破缺对称性的规范理论的对称结构在高能极限下得以恢复。

现在让我们把这些结果与弱相互作用规范理论中 $W^{+}W^{-}$ 产生振幅的直接计算相比较。从初始 $e^-_R$ 的情形开始。由于电子到 $W^-$ 的耦合是纯左旋的，图~\ref{fig:21.7} 的第三个图在这种情况下为零，所以计算稍微容易一些。图~\ref{fig:21.7} 的前两个图具有完全相同的结构，求和为
\begin{equation}
  i\mathcal{M}(e^-_Re^+_L \to W^+W^-) = \bar{v}_R\gamma^\mu u_R\ \frac{
    ie^2}{\sin^2\theta_W}\ \Bigl[ \cdots \Bigr]\
    \epsilon^{*}_\mu(k_+)\epsilon^{*}_\nu(k_-) ,
  \label{eq:21.96}
\end{equation}
其中方括号中的表达式包含光子和 $Z^0$ 传播子以及 $WW\gamma$ 和 $WWZ$ 顶角。这个方程在任一 $R_\xi$ 规范中都有效，因为如果我们忽略电子质量，
\begin{equation}
  q^\mu\ \bar{v}_R\gamma_\mu u_R = 0 .
  \label{eq:21.97}
\end{equation}
式~\eqref{eq:21.96} 的第二行包含上面提到的纵向 $W$ 玻色子的增强。如果我们用 \eqref{eq:21.60} 近似纵向极化矢量并丢掉不随 $s\to\infty$ 增长的项，这一行变得正比于
\begin{equation}
  \frac{1}{m_W^2}\ (k_+ - k_-)^\lambda\ \cdots ,
  \label{eq:21.98}
\end{equation}
它正是戈德斯通玻色子振幅 \eqref{eq:21.93} 的结构。进行代数运算，人们确认产生两个纵向 $W$ 玻色子的振幅等于产生两个戈德斯通玻色子的振幅，这与等价定理一致。

戈德斯通玻色子等价定理有重要的物理后果。在远大于 $m_W$ 的能量下，纵向 $W$ 和 $Z$ 玻色子的产生相对于朴素预期被增强，截面由戈德斯通玻色子相互作用支配。这意味着，例如，$e^{+}e^{-}\to W^{+}W^{-}$ 的截面随能量的增长方式由希格斯扇区的结构决定。在非常高的能量下，除非 $W$ 和 $Z$ 玻色子相互作用满足等价定理施加的关系，否则理论将破坏幺正性。这个要求是希格斯机制的原始动机之一。

\section{弱相互作用规范理论的一圈修正}

在第~11.6 节，我们给出了一个一般论证：在具有自发破缺对称性的理论中，经典级别成立的参数与可观测量之间的关系可以被圈修正所修改。这些修正将是有限的，并且事实上将是场论量子结构的确定预言。虽然我们最初是在具有自发破缺整体对称性的理论中讨论这个想法，但它同样适用于具有自发破缺规范对称性的理论。在本节中，我们将运用这个想法来推导弱相互作用规范理论中关系式的有限一圈修正。

在电弱理论中很容易找到零阶的自然关系。第~20.2 节给出的领头阶预言涉及相对较少的自由参数。这些预言中的许多是针对可以忽略夸克和轻子质量的能量做出的；于是它们只依赖于耦合常数 $g$ 和 $g'$ 以及真空期望值 $v$，后者设定了自发对称性破缺的标度。弱相互作用理论的其余成分用这些参数给出；例如，
\begin{equation}
  m_W = g\frac{v}{2},\qquad
  m_Z = \sqrt{g^2 + g'^2}\,\frac{v}{2},\qquad
  e = \frac{gg'}{\sqrt{g^2+g'^2}},\qquad
  G_F = \frac{1}{\sqrt{2}}\ \frac{g^2}{m_W^2} .
  \label{eq:21.109}
\end{equation}
即使在这组量中，我们也有四个依赖于三个基本参数的关系式，所以存在一个与拉格朗日量参数无关的可观测量关系式。

由于弱相互作用规范理论的许多预言由参数 $\sin^2\theta_W$ 决定，用可观测量来定义 $\sin^2\theta_W$ 然后把这个定义用作构造自然关系的基础是有用的。在第~20.2 节对电弱理论精确检验的讨论中，我们使用了定义
\begin{equation}
  \sin^2\theta^s_W = 1 - \frac{m_W^2}{m_Z^2} .
  \label{eq:21.110}
\end{equation}
作为比较不同实验的标准。但由于三个最精确已知的弱相互作用可观测量是 $\alpha$、$G_F$ 和 $m_Z$，基于这三个量构造 $\sin^2\theta_W$ 的另一个物理定义是有用的。定义 $\theta_0$ 使得
\begin{equation}
  \sin^2\theta_0\cos^2\theta_0 = \frac{\pi\alpha^{*}}{\sqrt{2}G_Fm_Z^2} ,
  \label{eq:21.111}
\end{equation}
其中 $\alpha^{*}$ 是在标度 $Q^2 = m_Z^2$ 处求值的 QED 跑动耦合常数。重整化群坚持认为，进入精确电弱预言的是弱相互作用标度处电荷的值，这一观察通过求和对涉及轻夸克和轻子的辐射修正图而被确认。式~\eqref{eq:21.111} 中诸量的当前最佳值给出
\begin{equation}
  s^2_0 = \sin^2\theta_0 = 0.2307 \pm 0.0005 .
  \label{eq:21.112}
\end{equation}
因此，这个量提供了一个非常精确的参考标准。

一旦取式~\eqref{eq:21.111} 定义 $\sin^2\theta_W$ 的参考值，第~20.2 节中把 $\sin^2\theta_W$ 与其他可观测量联系起来的方程就变成零阶自然关系。例如，树图方程
\begin{equation}
  \frac{m_W^2}{m_Z^2} = \cos^2\theta_W,\qquad
  \mathcal{A}_{LR} = \frac{\bigl( \tfrac{1}{2} - \sin^2\theta_W \bigr)^2
    - \bigl( \sin^2\theta_W \bigr)^2}
    {\bigl( \tfrac{1}{2} - \sin^2\theta_W \bigr)^2
    + \bigl( \sin^2\theta_W \bigr)^2} ,
  \label{eq:21.113}
\end{equation}
是联系弱相互作用四个可观测量的自然关系。这些关系式的修正将是理论的定义良好的预言。

原则上，我们现在可以计算修正参数 $m_Z$、$m_W$、$G_F$、$\alpha$ 和 $\mathcal{A}_{LR}$ 的所有一圈图。然而，这是一个非常复杂的练习，需要大量技术工具。%
\footnote{关于电弱理论一圈修正的详细理论讨论可以在 W.~Hollik，\emph{Fortschr.\ d.\ Physik} \textbf{38}，165 (1990) 中找到。} 在本节中，我们将集中讨论来自一个可以独立考虑的简单来源的辐射修正。除了反常问题之外，电弱理论不限制夸克或轻子代的数目。因此，计算一个夸克或轻子二重态贡献的一圈修正是合理且规范不变的。为明确起见，我们考虑 $(t, b)$ 夸克二重态的效应。

通过集中考虑重夸克引起的辐射修正，我们极大地简化了摆在面前的计算任务。弱相互作用规范理论的各种可观测量是从对初态和末态含轻费米子（轻子或夸克）的散射振幅的测量中提取的。例如，$G_F$ 从低能弱相互作用过程的强度测量，通常选择 $\mu$ 子衰变率：$\mu^{-}\to e^{-}\bar{\nu}_e\nu_\mu$。对任何这样的过程，都有多种一圈修正，如图~21.11 所示。除了对矢量玻色子传播子的修正之外，还有顶角修正、盒图以及带实光子发射的图。一般来说，各类图的贡献不是规范不变的；相反，规范不变性来自图~21.11(b)、(c) 和 (d) 各类图之间的抵消。然而，由于重夸克不直接耦合到轻轻子，$(t, b)$ 二重态只贡献图~21.11(f) 所示的单个图，它自身必然是规范不变的。同样的结论适用于对其他轻子弱相互作用过程的 $(t, b)$ 修正。如果我们忽略把 $t$ 和 $b$ 与其他物种混合的 CKM 角，这个结论也扩展到涉及轻夸克的弱相互作用过程。

对其他种类的粒子，例如希格斯扇区的粒子，也出现类似情形。希格斯扇区粒子到轻夸克或轻子的耦合正比于费米子的质量，我们常常可以忽略它。因此来自希格斯扇区粒子最重要的贡献是传播子修正。自发对称性破缺由单个标量场 $\phi$ 产生的情形特别容易分析；这将在习题~21.4 中完成。不直接耦合到外费米子的粒子的圈修正常被称为\emph{斜}修正，因为它们只是间接地进入低能弱相互作用。

\subsection{重夸克修正的影响}

那么，我们的任务是计算由 $(t, b)$ 二重态引起的关系式 \eqref{eq:21.113} 的修正。这两个关系式依赖于五个可观测量---$m_Z$、$m_W$、$\mathcal{A}_{LR}$、$\alpha$ 和 $G_F$---其中最后两个参数通过 $\theta_W$ 和式~\eqref{eq:21.111} 进入。我们将把这五个量表示为裸参数 $g$、$g'$ 和 $v$ 的函数，其修正正比于 $t$ 和 $b$ 真空极化图的组合。当我们计算关系式 \eqref{eq:21.113} 的修正时，零阶项将自然地消去。

我们所需的圈振幅是 $t$ 和 $b$ 夸克的一圈真空极化图。为了最直接地处理这些贡献，我们为真空极化振幅引入统一的记号。把涉及规范玻色子 $I$ 和 $J$ 的真空极化振幅记为
\begin{equation}
  i\Pi^{\mu\nu}_{IJ}(q) ,
  \label{eq:21.114}
\end{equation}
其中 $I$ 和 $J$ 可以是 $\gamma$、$W$ 或 $Z$。当规范玻色子有质量时，真空极化振幅自身不必是横向的，因此 $\Pi^{\mu\nu}_{IJ}(q)$ 不必在 $q^2 = 0$ 处为零。因此，我们将改变我们在 QED 情形中的记号，把 $\Pi^{\mu\nu}_{IJ}(q)$ 分解为张量结构的形式写成
\begin{equation}
  \Pi^{\mu\nu}_{IJ}(q) = \Pi_{IJ}(q^2)\,g^{\mu\nu} - \Delta_{IJ}(q^2)\,q^\mu q^\nu
  .
  \label{eq:21.115}
\end{equation}
在随后的所有例子中，因子 $q^\nu$ 将与轻轻子的流点乘，从而像式~\eqref{eq:21.97} 那样给出零。因此形状因子 $\Delta_{IJ}(q^2)$ 将从我们的计算中消去。我们先前关于 $\Pi^{\mu\nu}(q)$ 在 QED 中于 $q^2 = 0$ 处为零的结果在这个形式体系中表现为如下一组约束
\begin{equation}
  \Pi_{\gamma\gamma}(0) = \Pi_{\gamma Z}(0) = 0 .
  \label{eq:21.116}
\end{equation}
对其他振幅，我们的符号约定被选成使得 $\Pi_{IJ}(m^2)$ 的正值给规范玻色子一个正的质量移动。让我们还定义
\begin{equation}
  \Pi'_{IJ} = \frac{d\Pi_{IJ}}{dq^2}\Big|_{q^2=0},
  \label{eq:21.117}
\end{equation}
这就是我们在式~\eqref{eq:7.73} 中称为 $\Pi(0)$ 的量。

现在我们用这个记号写出对五个可观测量中每一个的圈修正。图~21.12 中的前两个图只是质量修正，因此直截了当地，
\begin{equation}
  m_Z^2 = (g^2 + g'^2)\frac{v^2}{4} + \Pi_{ZZ}(m_Z^2),\qquad
  m_W^2 = g^2\frac{v^2}{4} + \Pi_{WW}(m_W^2) .
  \label{eq:21.118}
\end{equation}
注意，两个真空极化振幅都在各自传播子的极点处求值。为了评估一圈修正引起的 $\alpha$ 移动，我们考虑图~21.12(c) 对低能 Coulomb 势的影响。领头阶传播子的值与一圈修正相结合给出因子
\begin{equation}
  \frac{e^2}{q^2}\ \Bigl( 1 + \Pi'_{\gamma\gamma}(q^2)\ \cdots \Bigr) ,
  \label{eq:21.119}
\end{equation}
其中在这个方程中，$e^2$ 像 \eqref{eq:21.109} 那样用裸变量给出。因此，在极限 $q^2\to 0$ 下，$\alpha$ 的观测值按照下式被修改
\begin{equation}
  \alpha = \alpha_0\ \Bigl( 1 - \Pi'_{\gamma\gamma}(0)\ \cdots \Bigr) .
  \label{eq:21.120}
\end{equation}
类似地，对 $G_F$ 和 $\mathcal{A}_{LR}$ 的修正可以用 $(t,b)$ 二重态的真空极化振幅来表达。

$(t, b)$ 二重态对真空极化振幅的贡献可以显式计算。这些图是圈中含 $t$ 和 $b$ 夸克的标准一圈真空极化图。结果是，对自然关系的修正依赖于 $t$ 和 $b$ 夸克质量平方之差，
\begin{equation}
  \Delta\rho \propto \frac{m_t^2 - m_b^2}{m_W^2} .
  \label{eq:21.121}
\end{equation}
这就是著名的结果：电弱辐射修正对 $t$ 夸克质量的依赖是二次的，而不是对数的。量 $\Delta\rho$ 修改 $m_W$ 与 $m_Z$ 之间的关系，使得
\begin{equation}
  \frac{m_W^2}{m_Z^2} = \cos^2\theta_W\ \Bigl( 1 - \frac{\Delta\rho}
    {\cos^2\theta_W} \Bigr) .
  \label{eq:21.122}
\end{equation}
弱相互作用可观测量对 $m_t^2$ 的依赖是这个理论的一个独特预言，它在顶夸克被发现之前被用来估计顶夸克质量。

\subsection{$m_t$ 的影响}

弱相互作用可观测量的一圈修正由对顶夸克质量的二次依赖支配。这种依赖源于 $(t, b)$ 二重态是分裂的这一事实，其中 $m_t \gg m_b$。大的 $t$ 质量的同位旋破缺效应修改了中性流与带电流相互作用的相对强度。

完整计算的结果是，自然关系 \eqref{eq:21.113} 被一个定义明确的量所修正，这个量依赖于 $m_t$ 和希格斯扇区的性质。把预言的修正与弱相互作用可观测量的实验值相比较，为这个理论提供了严格的检验，理论与实验的一致是标准模型的伟大成功之一。特别是，$\rho$ 参数的观测值与直接测量到的顶夸克质量范围以及轻希格斯玻色子一致。这种一致使人们对标准模型在量子层面的有效性有信心。

弱相互作用规范理论的辐射修正也为约束标准模型之外新粒子的性质提供了有力途径。任何耦合到 $W$ 和 $Z$ 玻色子的新粒子都将贡献于真空极化振幅，从而贡献于对自然关系的修正。理论与实验对已测量可观测量的精确一致对这类新粒子的质量和耦合施加了很强的约束。例如，这种方法已被用来约束希格斯玻色子的质量以及额外的 $Z'$ 玻色子可能存在的可能性。

\problems

\textbf{21.1} \emph{阿贝尔模型中的戈德斯通玻色子等价定理。} 在本习题中，验证第~21.1 节阿贝尔自发破缺规范理论中的戈德斯通玻色子等价定理。计算费米子-反费米子湮灭产生一个纵向规范玻色子的振幅，并证明它等于在高能下产生戈德斯通玻色子的振幅。

\textbf{21.2} \emph{$S$ 矩阵元的 $\xi$ 无关性。} 在第~21.1 节的阿贝尔模型中，通过对图~\ref{fig:21.2} 的图求和，显式验证树级费米子-费米子散射振幅与规范参数 $\xi$ 无关。

\textbf{21.3} \emph{$R_\xi$ 规范中的 $W$ 玻色子传播子。} 通过求矩阵 \eqref{eq:21.16} 的逆，验证规范玻色子传播子 \eqref{eq:21.17} 的形式。证明传播子中的非物理极点会在 $S$ 矩阵元中抵消。

\textbf{21.4} \emph{来自希格斯扇区的一圈修正。} 在阿贝尔自发破缺规范理论中，计算来自希格斯玻色子圈的、对规范玻色子质量以及对规范玻色子质量与耦合常数之间自然关系的一圈修正。证明这些修正是有限且定义良好的。
